\documentclass[12pt, a4paper]{article}

% Paquetes requeridos para el circuito
\usepackage[utf8]{inputenc}
\usepackage[T1]{fontenc}
\usepackage{circuitikz} 

\begin{document}

\begin{center}
    \begin{circuitikz}[american, scale=1.0, transform shape]
        
        % 1. Rama Izquierda: R3 (abajo) y R1 (arriba) en serie desde el nodo 'b' al 'a'
        \draw (0,0) to[R, l=$R_3$] (0,2)
                    to[R, l=$R_1$] (0,4);
        
        % 2. Rama Derecha: R4 (abajo) y R2 (arriba) en serie desde el nodo 'b' al 'a'
        % Se utiliza l_ para colocar las etiquetas de resistencia en la parte exterior
        \draw (4,0) to[R, l_=$R_4$] (4,2)
                    to[R, l_=$R_2$] (4,4);
        
        % 3. Rama Central: Fuente de tensión V conectada entre las dos uniones
        \draw (0,2) to[V, l=$V$] (4,2);
        
        % 4. Conexiones que cierran las ramas en paralelo (formando los raíles a y b)
        \draw (0,4) -- (4,4);
        \draw (0,0) -- (4,0);
        
        % 5. Terminales de salida explícitos 'a' y 'b'
        \draw (4,4) to[short, -o] (5.5,4) node[right] {a};
        \draw (4,0) to[short, -o] (5.5,0) node[right] {b};
        
    \end{circuitikz}
\end{center}

\end{document}